\documentclass[11pt,a4paper]{article}
\usepackage[T1]{fontenc}
\usepackage[utf8]{inputenc}
\usepackage{lmodern}
\usepackage[margin=25mm]{geometry}
\usepackage{amsmath,amssymb,amsthm}
\usepackage{xcolor}
\usepackage[hidelinks]{hyperref}
\newtheorem{theorem}{Theorem}
\newtheorem{proposition}[theorem]{Proposition}
\newtheorem{lemma}[theorem]{Lemma}
\newtheorem{corollary}[theorem]{Corollary}
\theoremstyle{plain}
\newtheorem{remark}[theorem]{Remark}
\newsavebox{\statementbox}
\newenvironment{statementframe}{%
  \par\addvspace{0.55\baselineskip}\noindent
  \begingroup\setlength{\fboxrule}{0.3pt}\setlength{\fboxsep}{7pt}%
  \begin{lrbox}{\statementbox}%
  \begin{minipage}{\dimexpr\linewidth-2\fboxsep-2\fboxrule\relax}%
}{%
  \end{minipage}\end{lrbox}%
  {\color{black!45}\fbox{\color{black}\usebox{\statementbox}}}%
  \endgroup\par\addvspace{0.45\baselineskip}%
}
\DeclareMathOperator{\Tr}{Tr}
\DeclareMathOperator{\diag}{diag}
\newcommand{\C}{\mathbb C}
\newcommand{\Q}{\mathbb Q}
\newcommand{\Z}{\mathbb Z}
\newcommand{\ket}[1]{\lvert#1\rangle}
\newcommand{\bra}[1]{\langle#1\rvert}
\newcommand{\repo}{https://github.com/jpbruneton/open_quantum_problems/blob/main/research/entries/complete-mub-entropy-ququart}
\setlength{\emergencystretch}{2em}
\setlength{\parindent}{0pt}
\setlength{\parskip}{0.3em plus 0.1em minus 0.1em}
\title{The five mutually unbiased bases of a ququart:\\
the sum of their Shannon entropies is at least seven bits}
\author{Opus 5.5\\[0.35em]
\normalsize Research led by Jean-Philippe Bruneton\\[-0.1em]
\small Universit\'e Paris Cit\'e, France}
\date{23 November 2026; Section~\ref{sec:d5two} added 27 November 2026}

\begin{document}
\maketitle
\begin{abstract}
Two orthonormal bases are mutually unbiased when a state that is certain in one
of them is completely random in the other. In a four-level system there are
five such bases, and quantum mechanics forbids the five measurement outcomes
from being simultaneously predictable. The entropic form of this
complementarity asks how small the sum of the five Shannon entropies can be.
Riccardi, Macchiavello and Maccone found the value of seven bits by numerical
minimization in 2017; the best proved lower bound was
$2+3\log_2 3\approx 6.755$ bits, due to S\'anchez-Ruiz. We prove that the sum
is at least $7$ bits for every state of $\C^4$ and every complete set of five
mutually unbiased bases, and that equality holds for exactly forty pure states.
In the two-qubit Pauli realization these are the stabilizer states whose
stabilizer group meets three of the five classes of commuting Pauli operators;
each gives the distribution $(\tfrac12,\tfrac12,0,0)$ in three bases and the
uniform distribution in the other two. The proof compares the Shannon entropy
with an explicit quartic symmetric polynomial and certifies two inequalities by
computer: a one-distribution inequality on the probability simplex, and a
degree-eight polynomial inequality on pure states. Both certificates are
replayed in exact rational or ball arithmetic. The same two-step scheme, with
an exact Hermitian sum-of-squares certificate, gives the non-tight bound
$10.2205$ bits for the six mutually unbiased bases of $\C^5$, above the
previously known $6\log_2 3\approx 9.510$ and below the numerical value
$10.2524$. We also prove that this numerical value is the exact minimum over
states supported on two vectors of one of the six bases, and that the $300$
pure states attaining it are strict local minima among all states.
\end{abstract}

\section{Introduction}

Mutually unbiased bases model maximally complementary measurements. Two
orthonormal bases $\{\ket{a}\}$ and $\{\ket{b}\}$ of $\C^d$ are mutually
unbiased when $|\langle a|b\rangle|^2=1/d$ for every pair of vectors, so that a
state prepared in a vector of one basis gives uniformly distributed outcomes in
the other. At most $d+1$ bases can be pairwise unbiased, and this number is
reached when $d$ is a prime power. Such a \emph{complete set} is
informationally complete: its outcome statistics determine the state.

The uncertainty principle for a complete set says that the $d+1$ outcome
distributions cannot all be sharp. Its entropic form bounds the sum of the
Shannon entropies from below. For a state $\rho$ and bases
$A_1,\dots,A_{d+1}$, write $p_{j,a}=\langle a_j|\rho|a_j\rangle$ for the
probability of outcome $a$ in basis $A_j$ and
\begin{equation}
  H(A_j)=-\sum_{a=1}^{d}p_{j,a}\log_2p_{j,a}
\end{equation}
for its Shannon entropy in bits, with $0\log_20=0$. The question is the exact
value of $\min_\rho\sum_jH(A_j)$ and the states that attain it. For a qubit the
answer is $2$ bits, attained by the eigenstates of the three Pauli
operators~\cite{WW2010,CBTW2017}. In dimension three the classical argument,
which uses only the purity of the state, gives $4$ bits, and this value is
attained~\cite{RMM2017}. From dimension four on, the classical bound is not
attained, and no tight value had been proved.

The four-level system, or ququart, is the first case where this argument
fails, and two qubits give the simplest picture of it. The fifteen non-trivial
two-qubit Pauli products split into five classes of three commuting operators,
\begin{equation}\label{eq:classes}
  \{XI,IX,XX\},\ \{ZI,IZ,ZZ\},\ \{YI,IY,YY\},\ \{XY,YZ,ZX\},\ \{XZ,YX,ZY\},
\end{equation}
and the common eigenbases of the five classes form a complete set of mutually
unbiased bases of $\C^4$. A common eigenvector of one class gives entropy $0$
in its own basis and $2$ bits in the four others, a total of $8$ bits. The
state $(\ket{00}+\ket{01})/\sqrt2$ does better. It is stabilized by $IX$, $ZI$
and $ZX$, which lie in three different classes, so it has the distribution
$(\tfrac12,\tfrac12,0,0)$, of entropy $1$, in three bases, and is uniform in the
other two. Its total is $3\cdot1+2\cdot2=7$ bits. Riccardi, Macchiavello and
Maccone found by numerical minimization that $7$ is the minimum~\cite{RMM2017},
and the value has since been tested experimentally~\cite{Serino2025}, but no
proof was known. The theorem below supplies it.

\begin{statementframe}
\begin{theorem}\label{thm:main}
Let $A_1,\dots,A_5$ be a complete set of mutually unbiased bases of $\C^4$.
For every density matrix $\rho$ on $\C^4$,
\begin{equation}\label{eq:main}
  \sum_{j=1}^{5}H(A_j)\ \ge\ 7 .
\end{equation}
Equality holds for exactly forty states, all pure. For the Pauli
realization~\eqref{eq:classes} they are the two-qubit stabilizer states whose
stabilizer group contains one operator from each of three different classes;
for each of the ten triples of classes there are four of them. Each minimizer
has the distribution $(\tfrac12,\tfrac12,0,0)$ in three bases and the uniform
distribution in the other two. For any other complete set, the minimizers are
the images of these forty states under the equivalence that maps the set onto
the Pauli realization.
\end{theorem}
\end{statementframe}

The forty minimizers can also be described without Pauli operators: they are
the states $(\ket{a}+e^{i\varphi}\ket{b})/\sqrt2$, with $a\ne b$ two vectors of
one basis, that have the distribution $(\tfrac12,\tfrac12,0,0)$ in two further
bases. They contain the twenty-four states of this form listed by Riccardi,
Macchiavello and Maccone for the first basis~\cite[Eq.~(52)]{RMM2017}. The
theorem therefore confirms their numerical value and shows that their list,
closed under the symmetries of the complete set, exhausts the minimizers.

The proof is computer-assisted. It compares $H$ with an explicit symmetric
polynomial $F$ of degree four in the probabilities and certifies two
inequalities: $H(p)\ge F(p)$ for every distribution on four outcomes, and
$\sum_jF(p_j)\ge7$ for every pure state. The first is a three-variable problem
on a simplex, the second a degree-eight polynomial inequality in six real
variables after a reduction by symmetries. Both are certified by
branch-and-bound searches whose leaves are replayed in exact rational
arithmetic or in ball arithmetic, together with exact local lemmas at the
equality points.

The same two-step scheme gives a new lower bound in dimension five. There the
classical bound $6\log_2 3\approx 9.510$ was the best one proved, while
numerical minimization gives $10.2524$~\cite{RMM2017}. Replacing the
branch-and-bound for the polynomial inequality by an exact Hermitian
sum-of-squares certificate, we obtain the following bound, which is not
tight.

\begin{statementframe}
\begin{theorem}\label{thm:d5}
For every complete set of six mutually unbiased bases $A_1,\dots,A_6$ of $\C^5$
and every density matrix $\rho$ on $\C^5$,
\begin{equation}
  \sum_{j=1}^{6}H(A_j)\ \ge\ \frac{20447}{2000}-\frac{6}{2000}=\frac{20441}{2000}=10.2205 .
\end{equation}
\end{theorem}
\end{statementframe}

\section{Related work}

\emph{Entropic uncertainty for complete sets.} The reviews of Wehner and
Winter~\cite{WW2010} and of Coles, Berta, Tomamichel and Wehner~\cite{CBTW2017}
describe the state of the art. For a complete set of mutually unbiased bases
and a pure state, the collision probabilities satisfy the identity
\begin{equation}\label{eq:collision}
  \sum_{j=1}^{d+1}\sum_{a=1}^{d}p_{j,a}^2=2,
\end{equation}
a consequence of the fact that a complete set forms a complex projective
$2$-design~\cite{BW2007,CBTW2017}. Since the Shannon entropy dominates the
collision entropy $-\log_2\sum_ap_a^2$, concavity of the logarithm gives the
bound $\sum_jH(A_j)\ge(d+1)\log_2\frac{d+1}{2}$, attributed to Ivanovi\'c and
to S\'anchez-Ruiz~\cite[Eq.~(6)]{WW2010}, \cite[Eq.~(81)]{CBTW2017}. For even
$d$, S\'anchez-Ruiz improved it to
$\frac d2\log_2\frac d2+(\frac d2+1)\log_2(\frac d2+1)$, which is tight for
$d=2$~\cite{WW2010}, \cite[Eq.~(86)]{CBTW2017}. For $d=4$ this gives
$2+3\log_2 3\approx 6.755$ bits, the best bound proved before this note. For
$d=5$ the first bound gives $6\log_23\approx 9.510$, as recalled by Wu, Yu and
M{\o}lmer~\cite[Eq.~(20)]{WYM2009}, who also derived bounds for subsets of a
complete set. Design-based and purity-based arguments of the same kind give
R\'enyi and collision-entropy bounds~\cite{KG2020,Howard2015,Frembs2026}; for the
two-qubit Pauli classes, Frembs, Natale, Wever and H\"ohn derive the
collision-entropy bound $5\log_2\frac52\approx6.610$ from a purity
invariant~\cite[Thm.~4]{Frembs2026}. Min-entropy bounds for arbitrary numbers
of mutually unbiased bases were obtained by Mandayam, Wehner and
Balachandran~\cite{MWB2010}, who record that in $d=4$ tight relations were
known for the min-entropy with two and four bases and for the collision
entropy with three and five bases~\cite[Sec.~I.C]{MWB2010}. Wehner and Winter
called tight uncertainty relations for all mutually unbiased bases ``an
interesting open question''~\cite[Sec.~III.B.2]{WW2010}, and Coles et al.\
list the strongest bound for a given number of bases as
open~\cite[App.~A]{CBTW2017}.

Section~\ref{sec:reduction} explains why the S\'anchez-Ruiz value is the
natural limit of these methods in $d=4$: it is the minimum of the sum of the
smallest entropies compatible with the individual collision probabilities under
the constraint~\eqref{eq:collision}, and the distributions that realize it,
three times $(\tfrac13,\tfrac13,\tfrac13,0)$ and twice $(\tfrac12,\tfrac12,0,0)$,
do not come from any quantum state. The proof below therefore has to use the
correlations between different bases.

\emph{Numerical values and experiments.} Riccardi, Macchiavello and
Maccone~\cite{RMM2017} minimized numerically, with Mathematica, the sum of
entropies for every number of mutually unbiased bases in dimensions three to
five. For the complete set in $d=3$ their value, $4$ bits, coincides with
Ivanovi\'c's bound and is attained, so it is proved; the values $7$ for $d=4$
and $10.25$ for $d=5$ come from numerical minimization only. Serino et al.\ tested all these values with
frequency-bin photonic states~\cite{Serino2025} and showed that in $d=5$ the
value for three bases depends on the choice of the triple~\cite{Serino2025b}. A
2026 review by Chesi and Maccone still presents the $d=4$ and $d=5$ values for
complete sets as numerical~\cite[Table~1]{CM2026}.

\emph{Other uncertainty relations for subsets.} Ballester and Wehner showed
that for $d=r^2$ there are $r+1$ mutually unbiased bases for which the trivial
pairwise bound is attained~\cite{BW2007}; Ambainis showed that for three bases
of the standard construction in prime dimension the average entropy can be
close to $\frac12\log_2 d$~\cite{Ambainis2009}. Mandayam, Bandyopadhyay, Grassl
and Wootters showed that any three of the five Pauli classes~\eqref{eq:classes}
determine a fourth maximal commuting set built from their operators, whose
common eigenstates give the distribution $(\tfrac12,\tfrac12,0,0)$ in each of
the three bases, so that the collision-entropy bound for three bases in $d=4$ is
tight~\cite[Thm.~1 and Thm.~6]{MBGW2014}. These common eigenstates are exactly
the minimizers of Theorem~\ref{thm:main}: the forty minimizers are the
eigenstates of the ten maximal commuting sets attached to the ten triples of
classes. Niekamp, Kleinmann and G\"uhne studied entropic uncertainty for
stabilizer bases and showed that the Maassen--Uffink bound is tight for any
pair of them~\cite{NKG2012}. In odd prime dimension, Howard found that among
equatorial states the magic states, not the stabilizer states, minimize the
total collision entropy over the non-computational bases~\cite{Howard2015}; for the Shannon entropy of the
complete set in $d=4$ the minimizers are stabilizer states. For odd prime
$d\ge5$, Li proved the sharp lower bound on the total support of a state in the
standard complete set and classified its extremizers; for $d=5$ they are the
superpositions of two vectors of the computational basis~\cite{Li2026}. In $d=5$ these two-vector
states also attain the numerical minimum $10.2524$ of the entropy sum (our
computation, Section~\ref{sec:d5}).

\emph{Classification of complete sets.} Brierley, Weigert and Bengtsson
classified all sets of mutually unbiased bases in dimensions two to
five~\cite{BWB2010}. In $d=4$ and in $d=5$ the complete set is unique up to an
equivalence generated by a global unitary, rephasing and permutation of the
vectors of each basis, permutation of the bases, and complex conjugation
\cite[Sec.~3, Sec.~4 and App.~A]{BWB2010}. Each of these operations transports
states and preserves $\sum_jH(A_j)$, so it suffices to prove the theorems for
one complete set in each dimension.

A search of the literature, of arXiv and of the public repository
\texttt{the-omega-\allowbreak institute/\allowbreak trureturing}, made on 22
and 23 November 2026, found no earlier proof of Theorem~\ref{thm:main} or of a
bound above $6\log_23$ for $d=5$. The closest item in that repository treats
the qubit case, in which the three Pauli entropies sum to at least two bits.
The search could not reach a citation index, so the papers citing~\cite{RMM2017}
were not all examined.

\section{Reduction to two inequalities}\label{sec:reduction}

\emph{Concavity.} The probabilities $p_{j,a}$ are linear in $\rho$ and $H$ is
strictly concave, so $\sum_jH(A_j)$ is concave in $\rho$ and its minimum is
attained at pure states. If a mixed state
$\rho=\sum_k\lambda_k\ket{\psi_k}\bra{\psi_k}$ with $\lambda_k>0$ attained the
minimum, every $\psi_k$ would be a minimizer and, by strict concavity, all
$\psi_k$ would have the same twenty probabilities $p_{j,a}$; since a complete
set is informationally complete, all $\psi_k$ would coincide. The minimizers are
therefore pure, and by the classification of complete sets it suffices to treat
the Pauli realization~\eqref{eq:classes}. The projector onto the common
eigenvector with eigenvalues $s_1,s_2\in\{\pm1\}$ of the class
$\{P_1,P_2,P_1P_2\}$ is $(I+s_1P_1)(I+s_2P_2)/4$, a matrix with entries in
$\Q(i)$.

\emph{Why purity is not enough.} For a distribution $p$ on four outcomes let
$c_k(p)=\sum_ap_a^k$. If one keeps only the identity~\eqref{eq:collision},
$\sum_jc_2(p_j)=2$, and bounds each entropy by the smallest entropy
$\Phi(c)$ of a distribution with collision probability $c$, the resulting
problem $\min\{\sum_j\Phi(c_j):\sum_jc_j=2\}$ has the value
$3\log_23+2\approx6.755$, attained by three distributions
$(\tfrac13,\tfrac13,\tfrac13,0)$ and two distributions
$(\tfrac12,\tfrac12,0,0)$. This is the S\'anchez-Ruiz bound. A comparison
function that depends on $c_2$ alone cannot do better, and the proof uses higher
power sums.

\emph{The comparison function.} Define the symmetric quartic polynomial
\begin{equation}\label{eq:F}
  F(p)=t_0+t_2c_2+t_3c_3+t_4c_4+t_{22}c_2^2,
\end{equation}
\begin{equation*}
  t_0=\tfrac{2793}{640},\quad t_2=-\tfrac{21657}{1600},\quad
  t_3=\tfrac{1841}{100},\quad t_4=-\tfrac{1081}{100},\quad t_{22}=\tfrac{61}{100}.
\end{equation*}
The coefficients $t_0$ and $t_2$ are fixed by $F(\tfrac12,\tfrac12,0,0)=1$ and
$F(\tfrac14,\tfrac14,\tfrac14,\tfrac14)=2$, so that $F$ agrees with $H$ at the
distributions of the minimizers. The other three were chosen by a semidefinite
programme that maximizes the curvature of inequality~(B) below at its zeros;
the proof uses only the rational values above.

\begin{proposition}[Inequality A]\label{prop:A}
For every probability vector $p$ on four outcomes, $H(p)\ge F(p)$, with
equality only at the uniform vector and at the six permutations of
$(\tfrac12,\tfrac12,0,0)$.
\end{proposition}

\begin{proposition}[Inequality B]\label{prop:B}
For every unit vector $\psi\in\C^4$,
$\sum_{j=1}^5F\bigl(p_j(\psi)\bigr)\ge7$, with equality only at the forty
states of Theorem~\ref{thm:main}.
\end{proposition}

\begin{proof}[Proof of Theorem~\ref{thm:main} from Propositions~\ref{prop:A} and~\ref{prop:B}]
For a pure state, $\sum_jH(A_j)\ge\sum_jF(p_j)\ge7$. Equality in the second
inequality forces $\psi$ to be one of the forty states, and these states
satisfy $\sum_jH(A_j)=3\cdot1+2\cdot2=7$. The reduction above extends the
bound to mixed states and to every complete set and shows that there are no
other minimizers. The forty states are listed with their distributions in
exact arithmetic in the verification data (\texttt{data/minimizers.json});
each of them has exactly three Pauli products with expectation $\pm1$, one in
each of three different classes, and each of the ten triples of classes occurs
four times. The twenty common eigenvectors of single classes give $8$ bits.
\end{proof}

\section{Certification of inequality A}\label{sec:A}

The domain is the three-dimensional simplex of probability vectors, in the
coordinates $(p_1,p_2,p_3)$. The function $H-F$ vanishes at the seven equality
points, so the proof combines two exact local lemmas with a branch-and-bound
search away from them.

\emph{Near the uniform vector $u$.} There $H-F$ vanishes with zero gradient,
and its Hessian $-\frac1{\ln2}\bigl(\diag(1/p_i)+J/p_4\bigr)-\nabla^2F$, where
$J$ is the all-ones matrix, is positive definite on the box
$|p_i-\tfrac14|\le0.0015$. An exact Taylor expansion of $\nabla^2F$ at $u$
bounds the variation of the polynomial part, and the smallest eigenvalue
$1.18$ of the midpoint matrix exceeds the Frobenius norm $0.90$ of the radius
matrix. The function $H-F$ is therefore strictly convex on the box and
positive except at $u$.

\emph{Near $e=(\tfrac12,\tfrac12,0,0)$.} Write $s=p_3+p_4$ and
$t=(p_1-p_2)/2$, and let $\eta(x)=-x\log_2x$ and $h$ be the binary entropy.
Subadditivity of $\eta$, the mean value theorem for $\eta$ and for $F$ along the
segment that moves the mass $s$ back to the first two coordinates give
\begin{equation}
  H(p)-F(p)\ge\varphi(t)+s\bigl(\log_2(1/s)-e_1-K\bigr),\qquad
  \varphi(t)=h(\tfrac12+t)-F(\tfrac12+t,\tfrac12-t,0,0),
\end{equation}
on the region $s\le0.02$, $|t|\le0.04$. Here $e_1=-0.2907$ bounds $\eta'$ from
above on $[0.45,0.54]$ and $K=5.2919$, computed with upward rounding, bounds
the corresponding directional derivative of $F$ on a box that contains the
segment. Since $2^{-(e_1+K)}=0.0312>0.02$, the second term is positive for
$0<s\le0.02$. The function $\varphi$ satisfies $\varphi(0)=\varphi'(0)=0$ and
$\varphi''\ge0.33$ on $[-0.04,0.04]$ by interval arithmetic, so $\varphi\ge0$.
Hence $H-F>0$ on the region except at $e$, and the symmetry of $H$ and $F$
under permutations transports the lemma to the six permutations of $e$.

\emph{Elsewhere.} On a simplex $S$ with vertices $v_0,\dots,v_3$ and longest
edge $\ell$, let $L$ be the affine interpolant of $F$. Since $H$ is concave,
$H-L$ is concave and attains its minimum on $S$ at a vertex. If
$\nabla^2F\succeq-mI$ on $S$ with $m\ge0$, then
$F-L\le\frac m2\sum_i\lambda_i|p-v_i|^2\le\frac m2\cdot\frac38\ell^2$ for the
barycentric coordinates $\lambda_i$ of $p$, hence
\begin{equation}
  \min_S(H-F)\ \ge\ \min_i(H-F)(v_i)-\tfrac{3}{16}\,m\,\ell^2 .
\end{equation}
The constant $m$ is obtained from an outward-rounded interval enclosure of
$\nabla^2F$ on the bounding box of $S$ and Gershgorin's theorem. Simplices
whose bound is not positive are bisected along their longest edge; simplices
inside the regions of the local lemmas are discarded. The search proves
$207\,716$ simplices and discards $1\,159$. Vertex values are computed in double
precision with an absolute margin of $10^{-12}$, and every proved simplex is
re-checked with vertex values computed in ball arithmetic at $128$ bits from
the exact rational vertices; no simplex is rejected. The smallest certified
lower bound is $2.2\times10^{-10}$.

\section{Certification of inequality B}\label{sec:B}

\emph{Symmetries.} The $24$ permutation matrices of the computational basis,
the $16$ diagonal matrices $\diag(1,i^a,i^b,i^c)$ with $a+b+c\equiv0\pmod4$,
and complex conjugation map the set of twenty projectors onto itself; this is
checked in exact arithmetic. Two orthogonal projectors of the set always belong
to the same basis, because projectors from different bases have overlap
$\tfrac14$, so these maps permute the bases and preserve $\sum_jF(p_j)$ and
$\sum_jH(A_j)$. Sorting the moduli of the components, fixing the global phase,
rotating the second and third components into the first quadrant and, when
needed, conjugating, every state becomes equivalent to $\psi\propto(1,w)$ with
$w=u+iv\in\C^3$ in
\begin{equation}
  D=\{\,|w_3|\le|w_2|\le|w_1|\le1,\ 0\le v_1\le u_1,\ u_2\ge0,\ v_2\ge0\,\}.
\end{equation}

\emph{The polynomial.} Homogenizing~\eqref{eq:F} with $N=|\psi|^2$, inequality
B becomes $g(u,v)\ge0$ on $D$, where
\begin{equation}
  g=\sum_{j=1}^5\Bigl(t_0N^4+t_2c_2(\tilde p_j)N^2+t_3c_3(\tilde p_j)N
  +t_4c_4(\tilde p_j)+t_{22}c_2(\tilde p_j)^2\Bigr)-7N^4,
\end{equation}
with $\tilde p_{j,a}=\langle\psi|\Pi_{j,a}|\psi\rangle$ the unnormalized
probabilities of $\psi=(1,u+iv)$. It is a polynomial of degree eight in six
real variables with $340$ rational coefficients. The forty minimizers give
$28$ points of the chart $w\in\{0,\pm1,\pm i\}^3$ (those whose first component
has maximal modulus), of which three lie in $D$: $w=(1,0,0)$, $(1,i,i)$ and
$(1,1,-1)$.

\emph{Local lemmas.} At each of the $28$ chart points $z_0$, exact rational
arithmetic gives $g(z_0)=0$, $\nabla g(z_0)=0$ and the expansion
$g(z_0+h)=h^{\mathsf T}Q_0h+\sum_{|\beta|\ge3}a_\beta h^\beta$. An exact
$LDL^{\mathsf T}$ factorization shows that $Q_0-\mu I$ is positive definite, and
$\mu>\sum_{k\ge3}A_kR^{k-2}$ with $A_k=\sum_{|\beta|=k}|a_\beta|$. Since
$|h^\beta|\le|h|^{|\beta|}$, this gives $g(z_0+h)>0$ for $0<|h|\le R$. The radii
are $R=2361/262144\approx9.0\times10^{-3}$ and, at the points of type
$(1,1,-1)$, $R\approx3.7\times10^{-3}$.

\emph{Branch-and-bound.} The search starts from the box
$[0,1]\times[0,1]\times[-1,1]\times[0,\tfrac34]\times[0,1]\times[-1,1]$ in the
variables $(u_1,u_2,u_3,v_1,v_2,v_3)$, which contains $D$ because $v_1\le u_1$
and $u_1^2+v_1^2\le1$ give $v_1\le1/\sqrt2$. For a box $c+[-r,r]$ it computes the
exact Taylor expansion $g(c+h)=\sum_\beta a_\beta h^\beta$, splits it into the
constant term $a_0$, the linear part, the quadratic part $h^{\mathsf T}Qh$ and
the tail $T=\sum_{|\beta|\ge3}|a_\beta|r^\beta$, and tests three lower bounds:
\begin{align}
  L_1&=a_0-\sum_{\beta\ne0}|a_\beta|r^\beta,\\
  L_3&=a_0-\sum_i|a_{e_i}|r_i+m\sum_ir_i^2-T,\qquad Q-mI\succ0,\ m\le0,\\
  L_2&=q(h^*)+\min_{|h_i|\le r_i}\nabla q(h^*)\cdot(h-h^*)-T,\qquad Q\succ0,
\end{align}
where $q$ is the quadratic Taylor polynomial and $h^*$ a point of the box; $L_2$
uses the convexity of $q$. A box is proved when one bound is positive,
discarded when it is disjoint from $D$ or contained in a ball of a local lemma,
and bisected otherwise. The search proves $52\,892$ boxes, discards $15\,560$
boxes disjoint from $D$ and $8$ boxes inside the balls, and runs in about one
minute.

\emph{Exact replay.} The search runs in double precision with explicit
rounding margins, but the proof does not depend on them. For each of the
$52\,892$ leaves, whose corners are dyadic numbers, an independent program
recomputes the Taylor expansion of $g$ at the centre exactly over $\Q$ and
verifies one of $L_1$, $L_3$, $L_2$ in exact rational arithmetic, with the
eigenvalue conditions checked by exact $LDL^{\mathsf T}$ factorizations. It
finds $2\,455$ leaves proved by $L_1$, $49\,078$ by $L_3$ and $1\,359$ by $L_2$,
and no failure; the smallest exact lower bound is $1.6\times10^{-5}$. A second
independent program rebuilds the bisection tree from the root box in exact
rational arithmetic and checks that every node is a recorded leaf, is disjoint
from $D$, or lies in a ball of a local lemma, the tests being exact; no node is
left uncovered. Together with the local lemmas, this proves inequality B with
no rounding argument. The zeros of $g$ in $D$ are thus exactly the three chart
points above, whose orbits under the symmetries are the forty states.

\section{Dimension five: a certified bound}\label{sec:d5}

In $d=5$ the complete set is unique up to equivalence~\cite{BWB2010}, and we use
$\{I,F_5,DF_5,\dots,D^4F_5\}$, where $F_5$ is the Fourier matrix and
$D=\diag(\omega^{k^2})_{k\in\Z_5}$ with $\omega=e^{2\pi i/5}$. Numerical
minimization from $1\,500$ random starts always returns
\begin{equation}
  V_5=6+5\,h\!\left(\tfrac{5+\sqrt5}{10}\right)\approx10.2524481255,
\end{equation}
attained by states $(\ket{j}+e^{i\varphi}\ket{k})/\sqrt2$ of two vectors of the
computational basis, with the distribution $(\tfrac12,\tfrac12,0,0,0)$ in that
basis and
$\bigl(\tfrac{5+\sqrt5}{20},\tfrac{5+\sqrt5}{20},\tfrac{5-\sqrt5}{20},\tfrac{5-\sqrt5}{20},0\bigr)$
in the five others. This value agrees with the $10.25$ of~\cite{RMM2017}; it is
not proved here.

The scheme of Sections~\ref{sec:A} and~\ref{sec:B} fails for a tight bound in
$d=5$. For every symmetric polynomial comparison function built from
$c_2,\dots,c_5$, the function $\sum_jF(p_j)-V_5$ is constant to $10^{-12}$
along a curve through the minimizer, because the sums $\sum_jc_k(p_j)$ stay
constant along it for $k\le5$. The zero of inequality B is then degenerate and
no local lemma of the kind used above applies. A non-tight bound is
nevertheless within reach.

\emph{Step A.} Take
$F=\sum_m\theta_m\prod_ic_{k_i}$ with the seven monomials $1$, $c_2$, $c_3$,
$c_4$, $c_2^2$, $c_5$, $c_3c_2$ and dyadic coefficients $\theta_m$ listed in
the verification data. A branch-and-bound on the ordered region
$p_1\ge\dots\ge p_5$ of the four-dimensional simplex, with the same simplex
bound as in Section~\ref{sec:A} (constant $\frac m2\cdot\frac25\ell^2$ for five
vertices) and vertex values re-checked in ball arithmetic, proves
$H(p)\ge F(p)-\eta$ with $\eta=\frac1{2000}$ ($6\,652\,089$ simplices, no
rejection).

\emph{Step B.} The Hermitian form
$f(\psi)=N^{18-5}\bigl(\sum_jF_{\rm hom}(p_j(\psi))-V_BN^{5}\bigr)$, with
$N=|\psi|^2$ and $V_B=\frac{20447}{2000}$, has a Gram matrix $G$ in the basis
of monomials $\psi^\gamma$, $|\gamma|=18$. The sums over the vector label and
over the power of $D$ are Gauss sums over $\Z_5$, so all roots of unity
cancel and $G$ has rational entries. It is block diagonal according to the
charge $\sum_kk\gamma_k\bmod5$, with five blocks of size $1\,463$. For each
block $A$ (after an exact diagonal congruence), a floating-point Cholesky
factor of $A-\tau I$ is rounded to a matrix $\tilde L$ with dyadic entries, and
the residual $E=A-\tau I-\tilde L\tilde L^{\mathsf T}$ is computed exactly in
integer arithmetic; the check $\|E\|_F<\tau$ with $\tau=2^{-20}$ shows that
$A=\tilde L\tilde L^{\mathsf T}+(\tau I+E)$ is positive semidefinite. Hence
$f\ge0$ and $\sum_jF(p_j)\ge V_B$ on unit vectors. The same check with
$V_B=10.23$ fails, as it should.

\emph{Conclusion.} Steps A and B give
$\sum_jH(A_j)\ge\sum_j(F(p_j)-\eta)\ge V_B-6\eta=10.2205$, which proves
Theorem~\ref{thm:d5} with the reduction of Section~\ref{sec:reduction}. The gap
$V_5-V_B$ between the numerical value and the best value certified at lift
degree $n$ was $0.149$, $0.091$, $0.071$, $0.051$ and $0.029$ for
$n=10,12,14,16,18$, a slow decrease consistent with the degenerate zero.

\subsection{Two-component states and local minimality}\label{sec:d5two}

Two analytic results, added on 27 November 2026, locate the candidate value
$V_5$ more precisely. The first proves that it is the exact minimum over all
states supported on two vectors of one of the six bases. The second proves
that the states attaining it are strict local minima among all density
matrices, including the directions that populate the other three components.
Write $S(\rho)$ for the sum of the six Shannon entropies in bits and
$h(x)=-x\log_2x-(1-x)\log_2(1-x)$.

\begin{statementframe}
\begin{theorem}\label{thm:d5two}
Let $\rho$ be a density matrix on $\C^5$ supported on the span of two vectors
of one of the six bases $I,F_5,DF_5,\dots,D^4F_5$. Then $S(\rho)\ge V_5$.
For the computational basis, equality holds exactly for the pure states
$(\ket j-\omega^s\ket k)/\sqrt2$ with $j<k$ and $s\in\Z_5$. This gives $50$
rays in each basis and $300$ distinct rays in total; each has exactly four
nonzero components in each of the five other bases.
\end{theorem}
\end{statementframe}

The affine relabellings $x\mapsto ux+c$ of $\Z_5$ with $u\ne0$, the linear
phases $\ket x\mapsto\omega^{bx}\ket x$, $D$ and $F_5$ permute the six bases.
For $F_5$, completing the square in the Fourier transform of a chirp
$\omega^{ax^2+bx}$ with $a\ne0$ gives another chirp, and the chirps with $a=0$
are exchanged with the computational basis. The statement and the counts can
therefore be transported from the computational basis to every basis.

\begin{proof}
Let $\psi=\sqrt p\,\ket j+\sqrt{1-p}\,e^{i\phi}\ket k$ with $j\ne k$, and set
$r=2\sqrt{p(1-p)}$. Since $k-j$ is invertible modulo $5$, the distribution of
$\psi$ in each of the five other bases is a permutation of
$p_l=\bigl(1+r\cos(\phi+2\pi l/5)\bigr)/5$, $l=0,\dots,4$, for a phase $\phi$
that depends on the basis only through a shift by a multiple of $2\pi/5$. At
fixed $r$, minimizing the entropy of this distribution amounts to maximizing
$\sum_lf(x_l)$ with $f(x)=(1+x)\ln(1+x)$ and $x_l=r\cos(\phi+2\pi l/5)$. The
series $f(x)=x+\sum_{k\ge2}(-1)^kx^k/[k(k-1)]$ converges absolutely on
$[-1,1]$. The power $\cos^k\theta$ is a combination with positive coefficients
of $\cos(m\theta)$ for $m\le k$ with $m\equiv k\pmod 2$, and the average over
$l$ keeps only the frequencies $m\in5\Z$, each contributing $\cos(m\phi)$. At
$\phi=\pi$ every such term equals $(-1)^m=(-1)^k$, so each term of
$\sum_k(-1)^k\sum_lx_l^k/[k(k-1)]$ takes its largest possible value
simultaneously. For $r>0$ the frequency $5$ in the term $k=5$ forces
$\phi\equiv\pi\pmod{2\pi/5}$ at any maximum.

At this phase, the entropy deficit relative to the uniform distribution in
each of the five bases is a series in $u=r^2$ with nonnegative coefficients
of $u^{k/2}$, $k\ge2$, hence a convex function of $u$. The entropy in the
initial basis, $h\bigl((1-\sqrt{1-u})/2\bigr)$, is strictly concave in $u$:
its derivative is $\operatorname{artanh}(z)/(2z\ln2)$ with $z=\sqrt{1-u}$, and
$\operatorname{artanh}(z)/z$ increases strictly with $z$. The total is
therefore strictly concave in $u\in[0,1]$, and its minimum lies at an
endpoint. At $u=0$ the state is a basis vector and $S=5\log_25\approx11.61$.
At $u=1$ the distribution in each of the five other bases consists of a zero
and twice each of $(5\pm\sqrt5)/20$, which gives $S=V_5<5\log_25$. Strict
concavity excludes the intermediate weights, and the phase condition gives
the $50$ rays. For a mixed state on the same two-dimensional span, concavity
of the entropy gives $S(\rho)\ge V_5$. Equality would force all the pure
components of $\rho$ to have the same distributions in all six bases; since
the six bases determine a density matrix, $\rho$ would be pure.
\end{proof}

\begin{statementframe}
\begin{proposition}\label{prop:d5local}
Let $P=\ket\psi\bra\psi$ be one of the $300$ rays of
Theorem~\ref{thm:d5two}. For every density matrix $\rho$ on $\C^5$, set
$t=1-\Tr(P\rho)$. If $0<t<1$, then
\[
  S(\rho)-V_5\ \ge\ t\log_2(1/t)-48t .
\]
In particular $S(\rho)>V_5$ whenever $0<t<2^{-48}$.
\end{proposition}
\end{statementframe}

\begin{proof}
By the symmetries above we may take $\psi=(\ket0-\ket1)/\sqrt2$. Its outcomes
of probability zero are $\ket2,\ket3,\ket4$ and, in the basis $D^aF_5$, the
vector $z_a$ with components $\omega^{a(x^2-x)}/\sqrt5$. By orthogonality of
characters, $\sum_az_az_a^*$ is block diagonal with all-ones blocks on the
fibres $\{0,1\}$, $\{2,4\}$, $\{3\}$ of $x\mapsto x^2-x$. Adding the three
coordinate projectors, the sum $Z$ of the projectors onto all zero-probability
outcomes has spectrum $\{0,1,2,2,3\}$ and kernel $\C\psi$. Hence the total
probability of these outcomes under $\rho$ is at least $t$, while each of them
has probability at most $t$, since each projector is orthogonal to $\psi$.
Their contribution to $S(\rho)$ is therefore at least $t\log_2(1/t)$.

For the outcomes of positive probability $p$ under $P$, write
$\eta(x)=-x\log_2x$ and $G=\sum\eta'(p)\,\Pi$, the sum running over their
projectors $\Pi$. The maps $\ket x\mapsto\omega^{x(x-1)}\ket x$ and
$\ket x\mapsto\ket{1-x}$ preserve the six bases and the ray of $\psi$, so they
commute with $G$. In the eigenspace of eigenvalue $1$ of the first map, the
odd subspace of the second is exactly $\C\psi$, so $G\psi$ is proportional to
$\psi$ and $G$ has no matrix elements between $\psi$ and its orthogonal
complement. Writing $X=\rho-P$, the linear term $\Tr(GX)$ is then at least
$-2\|G\|\,t$. Since every positive probability lies in $(1/8,1/2]$, we have
$|\eta'(p)|<2$ and $\|G\|\le12$.

For the quadratic term, the inequality $s\ln s-s+1\le(s-1)^2$ gives
$\eta(y)\ge\eta(p)+\eta'(p)(y-p)-(y-p)^2/(p\ln2)$ for every $y\ge0$. The six
spaces of traceless diagonal matrices in the six bases are mutually orthogonal
for the Hilbert--Schmidt product, because the bases are mutually unbiased, so
$\sum(y-p)^2\le\Tr X^2\le2t$, the sum running over all outcomes. As $p>1/8$
and $\ln2>2/3$, the quadratic loss is at most $24t$. Adding the three
estimates gives $S(\rho)-V_5\ge t\log_2(1/t)-24t-24t$.
\end{proof}

By concavity of the entropies the minimum of $S$ is attained on pure states.
Theorem~\ref{thm:d5two} then shows that a pure state with $S<V_5$, if one
exists, has at least three nonzero components in each of the six bases, and
Proposition~\ref{prop:d5local} shows that it stays away from the $300$ known
minimizers. The global value remains between $10.2205$ and $V_5$. The checker
\href{\repo/checks/two_sparse.py}{\texttt{two\_sparse.py}} confirms numerically
the $300$ rays, their value, the spectrum of $Z$, the alignment of $G\psi$ with
$\psi$ for each ray, and $S\ge V_5$ on $2\,000$ random two-component states;
the proofs above do not rely on it.

\section{Limitations}

Both theorems rest on computer-assisted proofs. Every step that decides the
result is checked in exact rational or integer arithmetic, or in ball
arithmetic with outward rounding: the values at the minimizers, the
symmetries, the local lemmas, the leaves and the coverage of the
branch-and-bound for inequality~B, the vertex values for inequality~A, and the
Gram certificates in dimension five. The interval enclosures of Hessians in
the simplex searches use outward-rounded floating-point arithmetic. The
coefficients of $F$ were found by optimization, which plays no role in the
proof.

Theorem~\ref{thm:main} answers the question for $d=4$ only. In $d=5$ the value
$V_5$ is proved only for two-component states and as a strict local minimum
(Section~\ref{sec:d5two}); the global minimum and the full set of minimizers
remain conjectural. The degenerate zero described in Section~\ref{sec:d5}
suggests that a tight proof needs invariants of degree at least six or a
certificate that allows exact zeros. Theorem~\ref{thm:d5} is $0.032$ bits
below $V_5$. The explicit radius $2^{-48}$ of
Proposition~\ref{prop:d5local} is far from optimal and is not used to cover a
neighbourhood in any certificate.

The result has not been independently refereed. An internal review in a
separate session reran every certificate, wrote the exact coverage check, and
found only formal defects, which were corrected before publication.

\section{Verification materials}\label{sec:verif}

All files are in the public entry directory
\href{https://github.com/jpbruneton/open_quantum_problems/tree/main/research/entries/complete-mub-entropy-ququart}{\texttt{research/entries/complete-mub-entropy-ququart}}.
For $d=4$, the shared routines are
\href{\repo/checks/common.py}{\texttt{common.py}} (exact projectors, the
coefficients of $F$, the polynomial $g$) and
\href{\repo/checks/ia.py}{\texttt{ia.py}} (outward-rounded interval
arithmetic). The checkers are
\href{\repo/checks/minimizers.py}{\texttt{minimizers.py}} (the forty minimizers
and the twenty basis states, exact),
\href{\repo/checks/symmetries.py}{\texttt{symmetries.py}} (the symmetry group,
exact), \href{\repo/checks/prove_A.py}{\texttt{prove\_A.py}} (inequality~A),
\href{\repo/checks/prove_B.py}{\texttt{prove\_B.py}} (local lemmas and
branch-and-bound for inequality~B),
\href{\repo/checks/replay_B_exact.py}{\texttt{replay\_B\_exact.py}} (exact
replay of the leaves) and
\href{\repo/checks/coverage_B_exact.py}{\texttt{coverage\_B\_exact.py}} (exact
coverage of the tree). They write the records
\href{\repo/data/minimizers.json}{\texttt{minimizers.json}},
\href{\repo/data/symmetries.json}{\texttt{symmetries.json}},
\href{\repo/data/A_certificate.json}{\texttt{A\_certificate.json}},
\href{\repo/data/B_certificate.json}{\texttt{B\_certificate.json}},
\href{\repo/data/B_leaves.npz}{\texttt{B\_leaves.npz}} (the leaves),
\href{\repo/data/B_replay.json}{\texttt{B\_replay.json}} and
\href{\repo/data/B_coverage.json}{\texttt{B\_coverage.json}}.
For $d=5$, \href{\repo/checks/gram5.py}{\texttt{gram5.py}} gives the exact Gram
matrices, \href{\repo/checks/check_gram5.py}{\texttt{check\_gram5.py}} compares
them with a direct computation,
\href{\repo/checks/params5.py}{\texttt{params5.py}} holds the coefficients,
\href{\repo/checks/cert_A5.py}{\texttt{cert\_A5.py}} and
\href{\repo/checks/cert_B5.py}{\texttt{cert\_B5.py}} prove steps A and B, and
the records are
\href{\repo/data/A5_certificate.json}{\texttt{A5\_certificate.json}},
\href{\repo/data/B5_certificate.json}{\texttt{B5\_certificate.json}} and
\href{\repo/data/B5_negative_test.json}{\texttt{B5\_negative\_test.json}}.
The numerical cross-check of Section~\ref{sec:d5two},
\href{\repo/checks/two_sparse.py}{\texttt{two\_sparse.py}}, writes
\href{\repo/data/two_sparse.json}{\texttt{two\_sparse.json}} in a few seconds.
The checks need Python~3 with \texttt{numpy}, \texttt{scipy}, \texttt{sympy}
and \texttt{python-flint}. On a laptop, the $d=4$ checks take about twenty
minutes in total, most of it for the exact replay, and the $d=5$ checks about
twenty minutes.

\begin{thebibliography}{99}

\bibitem{RMM2017}
A. Riccardi, C. Macchiavello and L. Maccone,
\emph{Tight entropic uncertainty relations for systems with dimension three to five},
Phys. Rev. A \textbf{95}, 032109 (2017);
\href{https://doi.org/10.1103/PhysRevA.95.032109}{doi:10.1103/PhysRevA.95.032109};
\href{https://arxiv.org/abs/1701.04304v1}{arXiv:1701.04304v1}.

\bibitem{Serino2025}
L. Serino, G. Chesi, B. Brecht, L. Maccone, C. Macchiavello and C. Silberhorn,
\emph{Experimental entropic uncertainty relations in dimensions three to five},
\href{https://arxiv.org/abs/2507.05025v1}{arXiv:2507.05025v1} (2025);
published in Phys. Rev. A \textbf{113}, 032420 (2026).

\bibitem{Serino2025b}
L. Serino, G. Chesi, B. Brecht, L. Maccone, C. Macchiavello and C. Silberhorn,
\emph{Complementarity-based complementarity: the choice of mutually unbiased
observables shapes quantum uncertainty relations},
Phys. Rev. Research \textbf{7}, 033152 (2025);
\href{https://doi.org/10.1103/v24q-sl6n}{doi:10.1103/v24q-sl6n};
\href{https://arxiv.org/abs/2406.11395v4}{arXiv:2406.11395v4}.

\bibitem{CM2026}
G. Chesi and L. Maccone,
\emph{Quantum uncertainty and entropy},
\href{https://arxiv.org/abs/2604.09384v1}{arXiv:2604.09384v1} (2026), book chapter,
\href{https://doi.org/10.1016/B978-0-443-33965-3.00002-X}{doi:10.1016/B978-0-443-33965-3.00002-X}.

\bibitem{WW2010}
S. Wehner and A. Winter,
\emph{Entropic uncertainty relations---a survey},
New J. Phys. \textbf{12}, 025009 (2010);
\href{https://doi.org/10.1088/1367-2630/12/2/025009}{doi:10.1088/1367-2630/12/2/025009};
\href{https://arxiv.org/abs/0907.3704v1}{arXiv:0907.3704v1}.

\bibitem{CBTW2017}
P. J. Coles, M. Berta, M. Tomamichel and S. Wehner,
\emph{Entropic uncertainty relations and their applications},
Rev. Mod. Phys. \textbf{89}, 015002 (2017);
\href{https://doi.org/10.1103/RevModPhys.89.015002}{doi:10.1103/RevModPhys.89.015002};
\href{https://arxiv.org/abs/1511.04857v2}{arXiv:1511.04857v2}.

\bibitem{WYM2009}
S. Wu, S. Yu and K. M{\o}lmer,
\emph{Entropic uncertainty relation for mutually unbiased bases},
Phys. Rev. A \textbf{79}, 022104 (2009);
\href{https://doi.org/10.1103/PhysRevA.79.022104}{doi:10.1103/PhysRevA.79.022104};
\href{https://arxiv.org/abs/0811.2298}{arXiv:0811.2298}.

\bibitem{BW2007}
M. A. Ballester and S. Wehner,
\emph{Entropic uncertainty relations and locking: tight bounds for mutually unbiased bases},
Phys. Rev. A \textbf{75}, 022319 (2007);
\href{https://doi.org/10.1103/PhysRevA.75.022319}{doi:10.1103/PhysRevA.75.022319};
\href{https://arxiv.org/abs/quant-ph/0606244v4}{arXiv:quant-ph/0606244v4}.

\bibitem{MWB2010}
P. Mandayam, S. Wehner and N. Balachandran,
\emph{A transform of complementary aspects with applications to entropic uncertainty relations},
J. Math. Phys. \textbf{51}, 082201 (2010);
\href{https://doi.org/10.1063/1.3477319}{doi:10.1063/1.3477319};
\href{https://arxiv.org/abs/1004.5086v2}{arXiv:1004.5086v2}.

\bibitem{MBGW2014}
P. Mandayam, S. Bandyopadhyay, M. Grassl and W. K. Wootters,
\emph{Unextendible mutually unbiased bases from Pauli classes},
Quantum Inf. Comput. \textbf{14}, 823--844 (2014);
\href{https://arxiv.org/abs/1302.3709v2}{arXiv:1302.3709v2}.

\bibitem{NKG2012}
S. Niekamp, M. Kleinmann and O. G\"uhne,
\emph{Entropic uncertainty relations and the stabilizer formalism},
J. Math. Phys. \textbf{53}, 012202 (2012);
\href{https://doi.org/10.1063/1.3678200}{doi:10.1063/1.3678200};
\href{https://arxiv.org/abs/1103.2316v2}{arXiv:1103.2316v2}.

\bibitem{Frembs2026}
M. Frembs, G. Natale, C. S. P. Wever and P. A. H\"ohn,
\emph{Maximal complementarity in the $n$-qubit Pauli group},
\href{https://arxiv.org/abs/2607.24988v1}{arXiv:2607.24988v1} (2026).

\bibitem{KG2020}
A. Ketterer and O. G\"uhne,
\emph{Entropic uncertainty relations from quantum designs},
Phys. Rev. Research \textbf{2}, 023130 (2020);
\href{https://doi.org/10.1103/PhysRevResearch.2.023130}{doi:10.1103/PhysRevResearch.2.023130};
\href{https://arxiv.org/abs/1911.07533v1}{arXiv:1911.07533v1}.

\bibitem{Howard2015}
M. Howard,
\emph{Maximum nonlocality and minimum uncertainty using magic states},
Phys. Rev. A \textbf{91}, 042103 (2015);
\href{https://doi.org/10.1103/PhysRevA.91.042103}{doi:10.1103/PhysRevA.91.042103};
\href{https://arxiv.org/abs/1501.05319v1}{arXiv:1501.05319v1}.

\bibitem{Ambainis2009}
A. Ambainis,
\emph{Limits on entropic uncertainty relations for 3 and more MUBs},
\href{https://arxiv.org/abs/0909.3720v2}{arXiv:0909.3720v2} (2009).

\bibitem{Li2026}
D. Li,
\emph{Sharp total-support uncertainty for the standard complete set of mutually unbiased bases},
\href{https://arxiv.org/abs/2609.16030v2}{arXiv:2609.16030v2} (2026).

\bibitem{BWB2010}
S. Brierley, S. Weigert and I. Bengtsson,
\emph{All mutually unbiased bases in dimensions two to five},
Quantum Inf. Comput. \textbf{10}, 803--820 (2010);
\href{https://arxiv.org/abs/0907.4097v2}{arXiv:0907.4097v2}.

\end{thebibliography}
\end{document}
